\section{Extension: elevated CO$_2$ and the quality of the nutrient supply}\label{sec:co2}
\paragraph{Purpose.} Sections~\ref{sec:nutrition} and \ref{sec:futurenut} count \emph{how much} of a crop is lost.
Elevated CO$_2$ also changes \emph{what is in} the grain: experiments find lower zinc, iron and protein in several C$_3$
crops. This section adds that channel as an \textbf{elevated-CO$_2$ experimental stress test} applied to the food supply
of 2021--23, keeps it separate from the yield channel, and combines the two by multiplication. It is not a calibrated
forecast for 2050. The index F4 is not altered.

\subsection{Evidence}
\paragraph{Which experiments.} \citet{myers2014} report meta-analysed changes in the edible portion of crops grown at
ambient and elevated CO$_2$ in seven FACE (free-air CO$_2$ enrichment) experiments in Japan, Australia and the USA,
with elevated CO$_2$ of 546--586\,ppm against ambient 363--386\,ppm. Their Extended Data Table~1 uses \emph{only these
experiments}; two later tables pool earlier FACE studies and then chamber studies, and change several values (wheat
zinc $-8.8$ and iron $-5.5\%$ when FACE data are pooled). The values in \file{code/config.py} match Extended Data Table~1
exactly, and that table is the evidence set used here. \citet{smith2018} draw on different inputs: zinc and iron from
the pooled FACE table, protein from \citet{medek2017}, an effect applied when $p<0.1$, and group averages for crops
without data (``C$_3$ legumes'', ``C$_3$ plants'', which include groundnut and rapeseed oils). The no-transfer rules
below exclude those group values. The paper's main text was not available, so only its abstract and supplementary
tables were read.

\paragraph{Parameter table.} Table~\ref{tab:co2params} gives the effects applied. The primary set is the original
FACE table for wheat, rice, maize and soybean (\cotwoNPrimary\ crop--nutrient effects, \cotwoNPrimarySig\ significant).
An extended variant adds \cotwoNExtended\ effects from other sources for chickpea, lentil and rapeseed. Full evidence,
with every candidate and the reason for rejection, is in \file{docs/co2\_nutrient\_effects.csv}; the selection with
reasons is \file{docs/co2\_effects\_used.csv}.

\begin{table}[H]\centering\scriptsize
\caption[CO$_2$ effects applied]{Elevated-CO$_2$ effects on the concentration in the edible part, as published. Primary
set: Extended Data Table~1 of \citet{myers2014}. The lentil values are derived from table means in the author's thesis
(wet and dry season are separate treatments, not a confidence interval); the published paper was read at abstract level.
Signif.: significance as reported.}\label{tab:co2params}
\adjustbox{max width=\linewidth}{\begin{tabular}{@{}llllllll@{}}
\toprule
Crop & Nutrient & Change \% (95\% CI) & Signif. & CO$_2$ elev. / amb. (ppm) & System & Source & Set \\
\midrule
Wheat & Iron & $-$5.1 ($-$6.5, $-$3.7) & yes & 546--550 / 382 & FACE & \citet{myers2014} & primary \\
Wheat & Protein & $-$6.3 ($-$7.5, $-$5.2) & yes & 546--550 / 382 & FACE & \citet{myers2014} & primary \\
Wheat & Zinc & $-$9.3 ($-$12.7, $-$5.9) & yes & 546--550 / 382 & FACE & \citet{myers2014} & primary \\
Rice & Iron & $-$5.2 ($-$7.6, $-$2.9) & yes & 570--576, 584 / 376--379, 386 & FACE & \citet{myers2014} & primary \\
Rice & Protein & $-$7.8 ($-$8.9, $-$6.8) & yes & 570--576, 584 / 376--379, 386 & FACE & \citet{myers2014} & primary \\
Rice & Zinc & $-$3.3 ($-$5.0, $-$1.7) & yes & 570--576, 584 / 376--379, 386 & FACE & \citet{myers2014} & primary \\
Maize & Iron & $-$5.8 ($-$10.9, $-$0.3) & yes & 550 / 385 & FACE & \citet{myers2014} & primary \\
Maize & Protein & $-$4.6 ($-$13.0, +4.5) & no & 550 / 385 & FACE & \citet{myers2014} & primary \\
Maize & Zinc & $-$5.2 ($-$10.7, +0.6) & no & 550 / 385 & FACE & \citet{myers2014} & primary \\
Soybean & Iron & $-$4.1 ($-$5.8, $-$2.5) & yes & 550 / 372--385 & FACE & \citet{myers2014} & primary \\
Soybean & Protein & +0.5 ($-$0.4, +1.3) & no & 550 / 372--385 & FACE & \citet{myers2014} & primary \\
Soybean & Zinc & $-$5.1 ($-$6.4, $-$3.9) & yes & 550 / 372--385 & FACE & \citet{myers2014} & primary \\
Chickpea & Protein & $-$13.5 ($-$21.4, $-$5.6) & yes & 500--700 / not reported & meta-analysis mix & \citet{medek2017} & extended \\
Rapeseed-mustard & Protein & $-$4.6 (CI not reported) & yes & 494 $\pm$ 16 / 388 $\pm$ 9 & FACE & \citet{hogy2010} & extended \\
Lentil & Iron & $-$4.4 (CI not reported) & yes & $\approx$550 / $\approx$400 & FACE & \citet{parvin2019} & extended (wet season) \\
Lentil & Zinc & $-$6.2 (CI not reported) & yes & $\approx$550 / $\approx$400 & FACE & \citet{parvin2019} & extended (wet season) \\
Lentil & Iron & $-$20.2 (CI not reported) & yes & $\approx$550 / $\approx$400 & FACE & \citet{parvin2019} & extended (dry season) \\
Lentil & Zinc & $-$27.3 (CI not reported) & yes & $\approx$550 / $\approx$400 & FACE & \citet{parvin2019} & extended (dry season) \\
\bottomrule
\end{tabular}
}
\end{table}

\paragraph{What is not estimated.} Table~\ref{tab:co2matrix} accounts for all \cotwoCellsTotal\ crop--nutrient
combinations: \cotwoCellsPrimary\ in the primary set, \cotwoCellsExt\ in the extended variant only, \cotwoCellsRev\ reviewed
but not used, and \cotwoCellsNE\ not estimated because no crop-specific evidence was found.
\begin{itemize}
\item \textbf{No transfer between species.} Field-pea effects are not applied to chickpea, lentil or bean, and soybean
effects are not applied to groundnut or rapeseed. Common bean and groundnut have no usable effect for any counted
nutrient. The only bean seed data are from a growth chamber at 800\,ppm; the two groundnut studies conflict.
\item \textbf{Energy and fat are not estimated for any crop.} No source reports a change in food energy, and fat or oil
data exist only for a few crops under single or conflicting studies.
\item \textbf{Extracted oils are untouched.} No source quantifies an effect on the energy or fat of extracted oil, so
seed concentration changes are never applied to the oil items.
\item \textbf{Reviewed but not used:} chickpea zinc (one open-top-chamber study, rescaled by a compilation assuming
linearity), rapeseed zinc (a value pooled over wheat, field pea and canola), and several ``no significant change, size
not reported'' results (chickpea and groundnut iron, rapeseed iron). The last group is recorded as \emph{no
detectable change}, not as zero.
\end{itemize}

\begin{table}[H]\centering\small
\caption[CO$_2$ evidence coverage]{Evidence for CO$_2$ effects on concentration. \textbf{S}: significant effect in the
primary set; \textbf{n.s.}: nonsignificant estimate in the primary set (applied with its interval); ext.: extended
variant only; rev.: evidence reviewed but not used; n/e: not estimated (no crop-specific evidence).}\label{tab:co2matrix}
\adjustbox{max width=\linewidth}{\begin{tabular}{@{}lccccc@{}}
\toprule
Crop & Energy & Protein & Fat & Zinc & Iron \\
\midrule
Wheat & n/e & \textbf{S} & n/e & \textbf{S} & \textbf{S} \\
Rice & n/e & \textbf{S} & n/e & \textbf{S} & \textbf{S} \\
Maize & n/e & \textbf{n.s.} & n/e & \textbf{n.s.} & \textbf{S} \\
Chickpea & n/e & ext. & n/e & rev. & rev. \\
Lentil & n/e & rev. & n/e & ext. & ext. \\
Common bean & n/e & rev. & rev. & rev. & rev. \\
Soybean & n/e & \textbf{n.s.} & rev. & \textbf{S} & \textbf{S} \\
Groundnut & n/e & rev. & rev. & rev. & rev. \\
Rapeseed-mustard & rev. & ext. & rev. & rev. & rev. \\
\bottomrule
\end{tabular}
}
\end{table}

\paragraph{Coverage of the supply.} Because only some crops and nutrients are covered, only part of each national
supply has a CO$_2$ effect applied (Table~\ref{tab:co2cover}). In India the covered crops supply \cotwoIndCovProtein\% of
protein, \cotwoIndCovZinc\% of zinc and \cotwoIndCovIron\% of iron, and none of the energy or fat. A full nine-crop CO$_2$
result is therefore unavailable. In the scenarios below, foods without an effect are held unchanged. This is a
\emph{scenario convention}, not evidence of no biological response.

\begin{table}[H]\centering\small
\caption[Supply covered by CO$_2$ effects]{Share of national supply to which a CO$_2$ effect is applied (left), and share
supplied by study-crop foods for which no effect could be estimated, the nutrient having no evidence for that crop (right).
Percent of the national total of the nutrient, 2021--23.}\label{tab:co2cover}
\adjustbox{max width=\linewidth}{\begin{tabular}{@{}lrrrrrrrr@{}}
\toprule
 & \multicolumn{5}{c}{Supply with a CO$_2$ effect applied (\%)} & \multicolumn{3}{c}{No evidence (\%)} \\
\midrule
 & Energy & Protein & Fat & Zinc & Iron & Fat & Zinc & Iron \\
India & 0 & 43 & 0 & 56 & 44 & 14 & 9 & 12 \\
World & 0 & 35 & 0 & 47 & 43 & 11 & 3 & 5 \\
Median of 20 countries & 0 & 30 & 0 & 45 & 41 & 8 & 3 & 6 \\
\bottomrule
\end{tabular}
}
\end{table}

\subsection{Can the experiment represent the future scenario?}
Not directly, so no scaling is done.
\begin{itemize}
\item The experiments contrast about 380 with about 550\,ppm. Under SSP2-4.5 the CO$_2$ level is 475, 507 and
537\,ppm in 2040, 2050 and 2060 \citep{ipcc2021annex3,meinshausen2020}, below the experiments' 546--586\,ppm, while the
2021--23 supply already reflects an observed global mean of about 417\,ppm (NOAA Global Monitoring Laboratory annual means, 2021--23). If the response were linear, the step from
today's level to 2050 would be about half the experimental step. Linearity is assumed in some compilations
\citep{terhaar2025}, but the primary sources give no dose--response: \citet{medek2017} found no reliable
dose-dependent decrease in protein. A remark about the direction of possible overstatement is therefore all that can be
said; nothing is calculated from it.
\item The effect is applied once, as a step change, to the whole eligible national supply of each covered crop,
domestic and imported, as measured in the edible part. Processing (milling, extraction) is not studied.
\item The FACE sites are in Australia, Japan and the USA; none is in India or the other tropical panel countries.
\end{itemize}
Reference levels with sources are in \file{docs/co2\_reference\_levels.csv}.

\subsection{Method}\label{sec:co2method}
Three scenarios are kept separate, all on the 2021--23 supply.
\begin{itemize}
\item \textbf{A. Quantity (yield) effects alone:} the conditional climate yield changes of Section~\ref{sec:futurenut},
losses only, applied to modelled crop--country pairs through $f=|\Delta Y|\,\text{SSR}^{*}$ (Eq.~\ref{eq:shock}). This is
a domestic production shock.
\item \textbf{B. Quality effects alone:} the concentration factor $q_{c,n}=1+p_{c,n}/100$ of Table~\ref{tab:co2params}
applied to the whole national supply of the covered crop's primary product, domestic and imported, in every country.
\item \textbf{C. Combined:} multiplicative accounting,
\begin{equation}\label{eq:co2comb}
S'_{c,n}=S_{c,n}\,\underbrace{(1-f_c)}_{\text{quantity left}}\;\underbrace{q_{c,n}}_{\text{concentration left}},
\end{equation}
never the sum of the two percentage changes. In India the sum of A and B for iron is \cotwoIndAddIron\% and the
combined result is \cotwoIndCombIron\%. The two differ wherever both act on the same food (independent recomputation: check R15).
\end{itemize}
Scopes are not mixed silently: A acts on domestic production of modelled pairs only, B acts on all eligible supply, and
C applies both to the same food item. Population growth, diet change, price and trade responses are not modelled.

\paragraph{Reference period.} The yield change of A and C is the \emph{hybrid counterfactual}: modelled climate change
from 1971--2000 to 2041--2060 applied to the 2021--23 supply. Because part of that change has already occurred, a
present-relative variant (A$'$, C$'$) measures the modelled 2041--2060 loss against the observed 2001--2024 level. Over
the 59 pairs the mean conditional yield change is \refMeanHybrid\% in the hybrid reference and \refMeanPresent\% in the
present-relative one; the difference (\refMeanObserved\%, about \refShareOccurred\% of the hybrid change) is the part
that the observed 2001--2024 climate already contains. Neither is a forecast of 2050 diets.

\paragraph{Envelopes.} The published 95\% limits are used as sensitivity envelopes, with every covered effect set to its
lower (or upper) limit at once. This is not a joint confidence interval: it assumes perfect correlation across crops and
ignores all other uncertainty in the model. Effects without a published interval are held at their central value in
the envelope.

\paragraph{Summary measure.} Nutrient outcomes stay visible. The combined summary is the equal-weighted mean percent
change of the three nutrients with evidence (protein, zinc and iron), the same equal-nutrient principle as MAR but
over three nutrients because energy and fat are not estimated. MAR itself is also reported; it is computed over all five
nutrients with energy and fat unchanged by convention, and its cap at 1 hides reductions that leave supply above the
requirement.

\subsection{Results}
Table~\ref{tab:co2india} and Figure~\ref{fig:co2} give India, and Table~\ref{tab:co2countries} the 20 panel countries.

\begin{table}[H]\centering\scriptsize
\caption[India: yield, quality and combined effects]{India: change in national supply (\%) by scenario, with the
sensitivity envelope [lower limit, upper limit] in brackets. A and A$'$: climate yield loss only; B: CO$_2$ quality only; C and
C$'$: combined. ``Mean of three'' is the equal-weighted mean of protein, zinc and iron. Iron NAR is the adequacy ratio
(baseline \cotwoIndIronNARzero).}\label{tab:co2india}
\adjustbox{max width=\linewidth}{\begin{tabular}{@{}lrrrrrr@{}}
\toprule
Scenario (India) & Protein \% & Zinc \% & Iron \% & Mean of three \% & $\Delta$MAR (pts) & Iron NAR \\
\midrule
A: climate yield loss only (hybrid) & $-$1.55 & $-$1.72 & $-$2.03 & $-$1.77 & $-$0.36 & 0.871 \\
B: CO$_2$ quality only & $-$2.96 [$-$3.56, $-$2.40] & $-$3.85 [$-$5.44, $-$2.27] & $-$2.24 [$-$3.04, $-$1.44] & $-$3.02 & $-$0.40 [$-$0.54, $-$0.26] & 0.869 \\
C: combined (hybrid) & $-$4.45 [$-$5.04, $-$3.90] & $-$5.44 [$-$6.99, $-$3.91] & $-$4.20 [$-$4.99, $-$3.42] & $-$4.70 & $-$0.75 [$-$0.89, $-$0.61] & 0.852 \\
A$'$: yield loss only (present-relative) & $-$0.78 & $-$0.81 & $-$1.03 & $-$0.88 & $-$0.18 & 0.880 \\
C$'$: combined (present-relative) & $-$3.72 [$-$4.31, $-$3.16] & $-$4.61 [$-$6.18, $-$3.05] & $-$3.25 [$-$4.04, $-$2.46] & $-$3.86 & $-$0.58 [$-$0.72, $-$0.44] & 0.860 \\
\bottomrule
\end{tabular}
}
\end{table}

\begin{figure}[H]\centering
\includegraphics[width=\linewidth]{X03_co2_quality.pdf}
\caption[CO$_2$ quality, yield and combined effects]{Left: India, change in national protein, zinc and iron supply under
climate yield loss (A), CO$_2$ concentration change (B) and both (C); whiskers are the published 95\% limits applied
together. Right: national zinc under the three scenarios for the 20 panel countries. Where A is zero, B and C coincide.}\label{fig:co2}
\end{figure}

\begin{itemize}
\item \textbf{India.} The quality channel alone lowers protein by \cotwoIndQualProteinAbs\%, zinc by \cotwoIndQualZincAbs\% and
iron by \cotwoIndQualIronAbs\%, which is \cotwoIndRatioProtein, \cotwoIndRatioZinc\ and \cotwoIndRatioIron\ times the climate
yield effect (A). Combined, supply falls by \cotwoIndCombProteinAbs\% (protein), \cotwoIndCombZincAbs\% (zinc) and
\cotwoIndCombIronAbs\% (iron); iron adequacy falls from \cotwoIndIronNARzero\ to \cotwoIndCombIronNAR\ and MAR by
\cotwoIndCombMARAbs\ points (envelope \cotwoIndCombHighMARAbs\ to \cotwoIndCombLowMARAbs). Against the present-relative reference, the
combined iron change is \cotwoIndCombPreIron\%.
\item \textbf{Quality exposes risk that yield alone misses.} In \cotwoNZeroQtyQual\ of the \cotwoNCountries\ countries the
climate scenario removes no modelled supply, yet the quality treatment lowers zinc by more than 1\%. In all
\cotwoNQualGtQtyZinc\ countries the quality effect on zinc exceeds the climate yield effect, and for iron in
\cotwoNQualGtQtyIron.
\item \textbf{Exposure follows dependence on the covered staples.} Zinc falls most in \cotwoTopZincCountries\ under the
combined scenario, countries with large wheat shares and, for Russia, a modelled wheat loss. This follows in part
arithmetically, since B is the supply-weighted mean of the effects, and says nothing about causes. The world's zinc,
iron and protein supply would fall by \cotwoWorldZincAbs, \cotwoWorldIronAbs\ and \cotwoWorldProteinAbs\% from quality alone.
\item \textbf{MAR hides most of it.} MAR falls by more than 0.05 points in only \cotwoNMarDrop\ countries
(\cotwoMarDropList), those where iron adequacy is already below 1. Zinc adequacy stays at or above \cotwoZincNARmin\ in every
country after the combined scenario, because the supply exceeds the requirement. Zinc and protein reductions are real
but are not shortfalls on this measure. No numbers of people with deficient intake are estimated, because that
requires population and intake distributions that were not used.
\item \textbf{Sensitivity.} Holding the three nonsignificant effects unchanged changes India's quality-only zinc to
\cotwoIndSigOnlyZinc\% and protein to \cotwoIndSigOnlyProtein\%. The extended variant with wet- or dry-season lentil
gives \cotwoIndExtWetZinc\% or \cotwoIndExtDryZinc\% for zinc in India; the dry-season lentil effect is several times the wet.
\end{itemize}

\begin{table}[H]\centering\scriptsize
\caption[National zinc, iron and protein changes]{National changes (\%) in the 20 panel countries. A: climate yield
loss only; B: CO$_2$ quality only; C: combined (hybrid reference, primary set). $\Delta$MAR in points.}\label{tab:co2countries}
\adjustbox{max width=\linewidth}{\begin{tabular}{@{}lrrrrrrrr@{}}
\toprule
Country & Zn A & Zn B & Zn C & Fe A & Fe B & Fe C & Protein C & $\Delta$MAR C \\
\midrule
Argentina & 0.00 & $-$3.63 & $-$3.63 & 0.00 & $-$2.43 & $-$2.43 & $-$1.64 & 0.00 \\
Australia & $-$0.09 & $-$2.73 & $-$2.82 & $-$0.13 & $-$1.70 & $-$1.83 & $-$1.25 & 0.00 \\
Bangladesh & $-$0.02 & $-$2.60 & $-$2.62 & $-$0.02 & $-$2.43 & $-$2.45 & $-$4.09 & $-$0.34 \\
Belarus & 0.00 & $-$2.35 & $-$2.35 & 0.00 & $-$1.36 & $-$1.36 & $-$0.79 & 0.00 \\
Brazil & $-$1.80 & $-$2.58 & $-$4.34 & $-$4.00 & $-$1.98 & $-$5.95 & $-$3.43 & $-$1.16 \\
Canada & $-$0.03 & $-$2.88 & $-$2.90 & $-$0.03 & $-$1.93 & $-$1.97 & $-$1.23 & 0.00 \\
China & $-$0.21 & $-$2.74 & $-$2.93 & $-$0.21 & $-$1.80 & $-$1.99 & $-$2.37 & 0.00 \\
Egypt & $-$0.15 & $-$5.59 & $-$5.72 & $-$0.14 & $-$3.52 & $-$3.66 & $-$3.52 & 0.00 \\
Ethiopia & 0.00 & $-$3.26 & $-$3.26 & 0.00 & $-$2.11 & $-$2.11 & $-$1.89 & 0.00 \\
France & 0.00 & $-$3.28 & $-$3.28 & 0.00 & $-$2.06 & $-$2.06 & $-$1.39 & 0.00 \\
India & $-$1.72 & $-$3.85 & $-$5.44 & $-$2.03 & $-$2.24 & $-$4.20 & $-$4.45 & $-$0.75 \\
Italy & $-$0.01 & $-$4.14 & $-$4.15 & 0.00 & $-$2.39 & $-$2.40 & $-$1.70 & 0.00 \\
Myanmar & $-$0.29 & $-$2.05 & $-$2.35 & $-$0.74 & $-$1.93 & $-$2.67 & $-$3.91 & $-$0.47 \\
Nigeria & 0.00 & $-$2.31 & $-$2.31 & 0.00 & $-$1.54 & $-$1.54 & $-$1.67 & $-$0.28 \\
Russia & $-$3.35 & $-$4.29 & $-$7.34 & $-$3.61 & $-$2.56 & $-$5.99 & $-$3.70 & 0.00 \\
South Africa & 0.00 & $-$3.99 & $-$3.99 & 0.00 & $-$3.41 & $-$3.41 & $-$2.42 & $-$0.56 \\
T\"urkiye & $-$0.03 & $-$4.51 & $-$4.54 & $-$0.04 & $-$2.38 & $-$2.41 & $-$2.31 & 0.00 \\
Ukraine & $-$1.08 & $-$4.02 & $-$5.05 & $-$1.07 & $-$2.33 & $-$3.33 & $-$2.28 & $-$0.02 \\
Tanzania & $-$1.57 & $-$2.75 & $-$4.32 & $-$3.47 & $-$2.41 & $-$5.88 & $-$4.86 & $-$0.98 \\
USA & $-$2.04 & $-$2.73 & $-$4.61 & $-$2.61 & $-$1.82 & $-$4.33 & $-$2.31 & 0.00 \\
\bottomrule
\end{tabular}
}
\end{table}

\subsection{Comparison with published studies}
Only conclusions and mechanisms are compared; our percentages are not set against theirs or against estimates of the
people affected.
\begin{itemize}
\item \textbf{Agreement in direction.} \citet{myers2014} and \citet{smith2018} find lower zinc, iron and protein in
C$_3$ staples under elevated CO$_2$ and conclude that diets relying on those staples are the most exposed;
\citet{zhu2018} (rice) and \citet{loladze2014} (a meta-analysis) also report lower concentrations of protein or minerals
in C$_3$ crops. Our results show the same direction and the same dependence on the covered staples. This is not
independent validation: the effect sizes are \emph{their} inputs.
\item \textbf{Crop and nutrient differences matter.} The effects are large and significant for wheat and rice, small
for soybean protein, and imprecise for maize (four pairs, one site). A single ``C$_3$ grain'' factor would hide this,
which is why group averages were not used.
\item \textbf{Where we differ.} \citet{smith2018} model population risk from intake distributions at 550\,ppm in 2050 and
fill missing crops with group averages; we estimate no populations, use no group averages, and show a supply-weighted
change for covered crops only. Their 2050 CO$_2$ level exceeds the 507\,ppm of SSP2-4.5. Our national energy and fat
are unchanged because no evidence supports a change.
\item \textbf{Yield and quality compound.} Calories and yield alone would miss the quality channel, but the combined
numbers also depend on the hybrid reference and on an experiment-based step that may overstate the SSP2-4.5 step.
\end{itemize}

\subsection{Limitations}
Evidence for common bean and groundnut is absent, for energy and fat absent, and for chickpea, lentil and rapeseed
thin; several sources are abstract-only. The response is assumed constant across cultivars, soils, nitrogen supply and
tropical conditions, although rice zinc responses are modified by cultivar and nitrogen \citep{myers2014}. National
nutrient supply is a mean, not intake. Elevated CO$_2$ would also raise yields of C$_3$ crops, which this extension does not
include.
